\section{Regularized logarithmic tests}\label{sec:tests}

Keep the profile parameters $a,\alpha,s_1,f,b,R_0$ from
Lemma~\ref{peaks:profiles}.  We shall construct a smooth
plurisubharmonic function at each degree whose average against its
representing density is uniformly negative.  The same negative bound
will hold against all later densities.  The estimates are uniform for
all sufficiently large integers $Q$, so the final value of $Q$ can be
chosen after the regularization and spacing parameters.

A positive measure representing holomorphic evaluation need not satisfy
the logarithmic inequality
$\log|f(0)|\le\int_S\log|f|\,d\mu$ for every holomorphic $f$
continuous on the closed ball. Hedenmalm gives an explicit example of
this distinction between representing and Jensen measures
\cite[Theorem~3.7]{Hedenmalm1989}.
Here we need smooth logarithmic tests and negative bounds that persist
through all later density corrections. We prove these properties
directly for the densities of Theorem~\ref{density:main}.

Fix a constant $C_P$ in the bound
$|P_k(z)|\le C_P|z|^k$ of Lemma~\ref{peaks:bounds}.

For $\epsilon>0$, define
\begin{equation}\label{tests:regularized}
 F_\epsilon(v)=
 \log\frac{|1-av|^2+\epsilon^2}{1+\epsilon^2},
 \qquad U_k(z)=F_\epsilon(P_k(z)).
\end{equation}
In particular $F_\epsilon(0)=U_k(0)=0$.  The function $U_k$ is smooth
and plurisubharmonic on $\mathbb C^2$, since
\begin{equation}\label{tests:hessian}
 (U_k)_{i\bar j}
 =\frac{a^2\epsilon^2}
        {(|1-aP_k|^2+\epsilon^2)^2}
   (P_k)_i\overline{(P_k)_j}.
\end{equation}
For fixed $\epsilon$ and $A<\infty$, all real derivatives of
$F_\epsilon$ of any fixed order are bounded on $|v|\le A$.
The mean-value estimate also gives
\begin{equation}\label{tests:lipschitz}
 |F_\epsilon(v)|\le C_{\epsilon,A}|v|
 \qquad (|v|\le A).
\end{equation}
All subsequent regularization constants may depend on the fixed
$\epsilon$.

\begin{lemma}[Exact averaging]\label{tests:averaging}
Let $k=k_j=Q^j$, where $Q$ satisfies the hypotheses of
Theorem~\ref{density:main}.  For $0\le r\le1$,
\begin{equation}\label{tests:averaging-identity}
 \int_S U_k(r\xi)W_j(\xi)\,d\sigma
 =
 \int_S U_k(r\xi)\,d\sigma
 +\sum_p\int_{\mathcal U_{k,p}}
     U_k(r\xi)e^{2y/k}f(y)\cos(k\theta)\,d\sigma.
\end{equation}
\end{lemma}

\begin{proof}
Under the global phase rotation $\xi\mapsto e^{i\omega}\xi$,
homogeneity gives $P_k(e^{i\omega}\xi)=e^{ik\omega}P_k(\xi)$.
Thus $U_k(r\xi)$ has only phase frequencies in $k\mathbb Z$.
By Lemma~\ref{density:moments}, the phase bandwidth of $W_{j-1}$
is $\ell_{j-1}<k$ and its constant coefficient on each orbit is one.
Averaging the product over the phase therefore gives
\[
 \int_S U_k(r\xi)W_{j-1}(\xi)\,d\sigma
 =\int_S U_k(r\xi)\,d\sigma.
\]

Each patch is phase-invariant.  Its radial variable $y$ is fixed
under rotation, while $\theta$ and $\gamma$ both increase by $\omega$.
Write the central restriction as
$S_p(\theta)=W_{j-1}(e^{i\theta}p)
 =\sum_s\widehat S_p(s)e^{is\theta}$.
In the patch increment of~\eqref{density:recursion}, the frequencies
are $\pm(k+s)$ with $|s|\le\ell_{j-1}<k$.  Pairing with a frequency in
$k\mathbb Z$ is possible only for $s=0$.
Since $\widehat S_p(0)=1$ and $S_p'$ has zero constant coefficient,
the surviving part of the increment
is exactly $e^{2y/k}f(y)\cos(k\theta)$.  Integrating this phase average
proves~\eqref{tests:averaging-identity}.
\end{proof}

The next estimate chooses the regularization before the spacing.
This order matters: the single-peak approximation will be applied
to a fixed smooth logarithm, and its error can then be made small by
increasing $D_0$.

\begin{proposition}[Negative averages at one scale]\label{tests:negative}
There exist $\epsilon>0$, $\eta\in(0,1/2)$ and $D_*\ge1$, depending
only on the fixed profile data, with the following property.
For every $D_0\ge D_*$ there are $k_*>0$ and $d_2>0$ such that, for
every integer $Q$ satisfying the density theorem and $Q\ge k_*$,
\begin{equation}\label{tests:one-scale-negative}
 \int_S U_{k_j}(r\xi)W_j(\xi)\,d\sigma\le-d_2
 \qquad
 \bigl(j\ge1,\quad 1-\eta\le r^{k_j}\le1\bigr).
\end{equation}
The constants $k_*,d_2$ may depend on $D_0$, but are independent of
$j$ and of further increases in $Q$.
\end{proposition}

\begin{proof}
Put
\[
 \Delta=2s_1(\cosh\alpha-1)-\alpha^2>0.
\]
We first compute the integral for a single peak:
\begin{equation}\label{tests:model}
 I_\epsilon(\rho)=
 \int_0^{R_0}\int_0^{2\pi}
 \log\frac{|1-a\rho e^{-y}e^{it'}|^2+\epsilon^2}
               {1+\epsilon^2}
 (1+f(y)\cos t')\,\frac{dt'}{2\pi}\,dy.
\end{equation}
At $\rho=1$, the unregularized logarithm has only one singularity,
at $(y,t')=(\alpha,0)$ modulo $2\pi$.  Near that point,
\[
 |1-e^{\alpha-y+it'}|^2
 \asymp (y-\alpha)^2+(t')^2.
\]
Its logarithm is locally integrable.  For $0<\epsilon\le1$, the
regularized logarithm is bounded above uniformly, and its negative
part is bounded by the negative part of the unregularized logarithm
plus $\log2$.  Dominated convergence therefore applies in
\eqref{tests:model}.

For $q<1$, the mean and cosine moment of
$\log|1-qe^{it'}|^2$ are $0$ and $-q$, respectively, by the logarithm
series.  For $q>1$, factoring out $q^2$ gives mean $2\log q$ and
cosine moment $-q^{-1}$.  Using the cancellation
$\int_0^{R_0}e^{-y}f(y)\,dy=0$, we obtain
\begin{align}
 I_0(1)
 &=\alpha^2-\int_0^\alpha f(y)e^{y-\alpha}\,dy
             -\int_\alpha^{R_0}f(y)e^{\alpha-y}\,dy \notag\\
 &=\alpha^2+\int_0^\alpha
       f(y)\bigl(e^{\alpha-y}-e^{-(\alpha-y)}\bigr)\,dy \notag\\
 &=\alpha^2-2s_1(\cosh\alpha-1)=-\Delta.
 \label{tests:model-sign}
\end{align}
The last equality uses $f=-s_1$ on $[0,\alpha]$.

We now compare this model with~\eqref{tests:averaging-identity}.
For $k=k_j$ put $\rho=r^k$, so homogeneity gives
$U_k(r\xi)=F_\epsilon(\rho P_k(\xi))$.
Off the patches, $|aP_k|\le1/2$ by Lemma~\ref{peaks:profiles}.
If $v=\rho P_k(\xi)$, $0\le\rho\le1$, and $s=|1-av|^2$, then
$1/4\le s\le9/4$ and
\[
 F_\epsilon(v)-\log|1-av|^2
 =\log(1+\epsilon^2/s)-\log(1+\epsilon^2).
\]
The derivative in $s$ of the first term has magnitude at most
$16\epsilon^2$, and $|s-1|\le C|v|$.  Hence
\[
 |F_\epsilon(v)-\log|1-av|^2|
 \le C\epsilon^2|P_k(\xi)|.
\]
The off-patch set is phase-invariant, and the phase mean of the
unregularized logarithm is zero.  By the $L^1$ estimate of
Lemma~\ref{peaks:bounds}, its total contribution is thus at most
\begin{equation}\label{tests:off-patch}
 C_{\rm off}\epsilon^2\,\frac{N_k}{k}.
\end{equation}
Here $C_{\rm off}$ is independent of $D_0,k,\rho$ and
$0<\epsilon\le1$.

On a patch, the measure formula of Lemma~\ref{peaks:patches} writes
the combined contribution as
\[
 \frac2k\int_0^{R_0}\int_{\mathbb T^2}
 F_\epsilon(\rho P_k(\xi))
       \bigl(e^{-2y/k}+f(y)\cos(k\theta)\bigr)
       \,\frac{d\theta\,d\gamma}{(2\pi)^2}\,dy.
\]
For fixed $\epsilon,R_0$, replace $P_k(\xi)$ by
$e^{-y}e^{ik\theta}$ using the single-peak estimate, and replace
$e^{-2y/k}$ by one.  The bounded derivative of $F_\epsilon$ and the
bound $|e^{-2y/k}-1|\le2R_0/k$ show that the result differs from
$(2/k)I_\epsilon(\rho)$ by at most
\begin{equation}\label{tests:patch-error}
 \frac2k C_{\epsilon,R_0}
 \left(e^{-c(D_0-\sqrt{2R_0})^2}+\frac1k\right).
\end{equation}
All constants are uniform over patches and $0\le\rho\le1$.
The change from $k\theta$ to $t'$ in the angular average is exact,
as $k$ is an integer.

Choose $\epsilon\in(0,1]$ so small that
\[
 |I_\epsilon(1)+\Delta|\le\Delta/8,
 \qquad C_{\rm off}\epsilon^2\le\Delta/8.
\]
By continuity of~\eqref{tests:model} in $\rho$, choose
$\eta\in(0,1/2)$ such that
$I_\epsilon(\rho)\le-3\Delta/4$ for $1-\eta\le\rho\le1$.
Next choose $D_*$ large enough for patch disjointness and for the
first error in~\eqref{tests:patch-error} to be at most $\Delta/16$
inside its factor $2/k$.  For each $D_0\ge D_*$, take $k_*$ large
enough for all the geometric estimates and for the remaining error
inside this factor to be at most $\Delta/16$.
Summing the patches and adding~\eqref{tests:off-patch} gives, in
particular,
\[
 \int_S U_k(r\xi)W_j\,d\sigma
 \le-\frac{\Delta}{2}\frac{N_k}{k}
 \le-\frac{\Delta}{2D_0^2}.
\]
Thus $d_2=\Delta/(2D_0^2)$ is a valid choice.
\end{proof}

\begin{lemma}[Negative averages against later densities]\label{tests:later}
Fix $\epsilon,\eta,D_0,k_*,d_2$ as in
Proposition~\ref{tests:negative}, and put $d_3=d_2/2$.
There is a threshold $Q_*$ such that, for every integer $Q\ge Q_*$,
the densities satisfy
\begin{equation}\label{tests:later-negative}
 \int_S U_{k_j}(r\xi)W_N(\xi)\,d\sigma\le-d_3
 \qquad
 \bigl(N\ge j\ge1,\quad1-\eta\le r^{k_j}\le1\bigr).
\end{equation}
\end{lemma}

\begin{proof}
Include the density and degree thresholds in $Q_*$, and require
$Q\ge4$.  By Lemma~\ref{density:moments}, the increment
$W_\nu-W_{\nu-1}$ has frequencies within $\ell_{\nu-1}$ of
$\pm k_\nu$.  Since $\ell_{\nu-1}\le2k_{\nu-1}$, each such
frequency has absolute value at least $(Q-2)k_{\nu-1}$.
Summing the increments with $\nu>j$ shows that every frequency in
$W_N-W_j$ has absolute value at least $(Q-2)k_j$.

For fixed $\xi,r$, write the logarithmic profile as
\[
 V(t')=\log\frac{|1-Ae^{it'}|^2+\epsilon^2}{1+\epsilon^2},
 \qquad A=a r^{k_j}P_{k_j}(\xi).
\]
The uniform polynomial bound gives $|A|\le aC_P$.  Differentiating
the logarithm shows that
\[
 \|V''\|_\infty\le
 B_\epsilon:=
 \frac{2aC_P}{\epsilon^2}
 +\frac{4a^2C_P^2}{\epsilon^4}.
\]
Two integrations by parts bound its $m$th Fourier coefficient by
$B_\epsilon/m^2$ for $m\ne0$. To track these coefficients as functions
of $\xi$, define
\[
 c_m(\xi)=\int_0^{2\pi}
 F_\epsilon\bigl(r^{k_j}P_{k_j}(\xi)e^{it'}\bigr)
 e^{-imt'}\,\frac{dt'}{2\pi}.
\]
Homogeneity and a change of variable give
$c_m(e^{i\omega}\xi)=e^{imk_j\omega}c_m(\xi)$.
Thus $c_m$ has global phase frequency $mk_j$, and
$U_{k_j}(r\xi)=\sum_{m\in\mathbb Z}c_m(\xi)$.
The coefficient bound makes this series uniformly convergent.
Its truncation to $|m|\le Q-3$ has uniform error at most
$2B_\epsilon/(Q-3)$ and only frequencies of absolute value less
than $(Q-2)k_j$. Phase averaging therefore makes its integral
against $W_N-W_j$ vanish.
Both densities are positive probabilities by
Theorem~\ref{density:main}; hence
$\|W_N-W_j\|_{L^1(\sigma)}\le2$.  It follows that
\[
 \left|\int_S U_{k_j}(r\xi)(W_N-W_j)\,d\sigma\right|
 \le\frac{4B_\epsilon}{Q-3}.
\]
Choose $Q_*$ so that this is at most $d_2/2$, and apply
\eqref{tests:one-scale-negative}.
\end{proof}
